%%
%% This is file `sample-sigconf.tex',
%% generated with the docstrip utility.
%%
%% The original source files were:
%%
%% samples.dtx  (with options: `all,proceedings,bibtex,sigconf')
%% 
%% IMPORTANT NOTICE:
%% 
%% For the copyright see the source file.
%% 
%% Any modified versions of this file must be renamed
%% with new filenames distinct from sample-sigconf.tex.
%% 
%% For distribution of the original source see the terms
%% for copying and modification in the file samples.dtx.
%% 
%% This generated file may be distributed as long as the
%% original source files, as listed above, are part of the
%% same distribution. (The sources need not necessarily be
%% in the same archive or directory.)
%%
\documentclass[manuscript,screen,nonacm]{acmart}
\usepackage{tabularx}
\usepackage{array}
\usepackage{booktabs}
\usepackage{colortbl}
\usepackage{xcolor}
\definecolor{trizgreen}{RGB}{0,125,70}
\definecolor{trizmagenta}{RGB}{160,45,120}
\definecolor{trizblue}{RGB}{70,110,140}
\definecolor{triznavy}{RGB}{10,50,75}
\definecolor{trizrow}{RGB}{240,244,247}

\AtBeginDocument{%
  \providecommand\BibTeX{{%
    Bib\TeX}}}

\setcopyright{none}
\settopmatter{printacmref=false}
\copyrightyear{2026}
\acmYear{2026}
\acmDOI{}
\acmISBN{}
\acmConference[AI4Science]
{AI4Science}
{2026}{Shenzhen, China}
\acmBooktitle{AI4Science, 2026, Shenzhen, China}
\settopmatter{printfolios=true}


\begin{document}
\newcommand{\triza}{TRIZ$^a$}
\newcommand{\papertitle}{\triza: Guiding Agent Evolution from Pattern Recognition to Solution Invention}
\newcommand{\complete}{\textcolor{trizgreen}{\textbf{complete}}}
\newcommand{\withheld}{\textcolor{trizmagenta}{\textbf{not promoted}}}

\title[TRIZa Supplementary Appendix]{Supplementary Appendix to \papertitle}


\author{Wenyin Liu}
\authornote{Corresponding author.}
\email{liuwy@gdut.edu.cn}
\affiliation{
  \institution{Guangdong University of Technology}
  \city{Guangzhou}
  \country{China}
}

\author{Yiheng Huang}
\email{huangyiheng.bupt@gmail.com}
\affiliation{
  \institution{Beijing University of Posts and Telecommunications}
  \city{Beijing}
  \country{China}
}

\author{Kai Wang}
\email{kaiwen@denglu.net.cn}
\affiliation{
  \institution{Beijing Denglu Technology Ltd}
  \city{Beijing}
  \country{China}
}

\renewcommand{\shortauthors}{Liu et al.}

\begin{abstract}
This independently readable supplement accompanies the eight-page AI4Science
main paper. It preserves the full TRIZ mapping, evaluation contracts,
qualification and controller rules, traceability material, and expanded
method and comparison descriptions condensed in the main paper.
No new model was trained or search run to prepare it. Original and public
redacted file identities are distinguished, and retained decision replay
is not treated as independent refitting from missing sample-level predictions.
The shared public entry point is https://triza.cn/triza-paper-exp-data2026/.
\end{abstract}

\begin{CCSXML}
<ccs2012>
<concept>
<concept_id>10010147.10010178.10010219.10010221</concept_id>
<concept_desc>Computing methodologies~Intelligent agents</concept_desc>
<concept_significance>500</concept_significance>
</concept>
<concept>
<concept_id>10010147.10010178.10010187</concept_id>
<concept_desc>Computing methodologies~Knowledge representation and reasoning</concept_desc>
<concept_significance>300</concept_significance>
</concept>
<concept>
<concept_id>10010147.10010257</concept_id>
<concept_desc>Computing methodologies~Machine learning</concept_desc>
<concept_significance>300</concept_significance>
</concept>
<concept>
<concept_id>10002978.10003029.10011703</concept_id>
<concept_desc>Security and privacy~Intrusion/anomaly detection and malware mitigation</concept_desc>
<concept_significance>300</concept_significance>
</concept>
</ccs2012>
\end{CCSXML}

\ccsdesc[500]{Computing methodologies~Intelligent agents}
\ccsdesc[300]{Computing methodologies~Knowledge representation and reasoning}
\ccsdesc[300]{Computing methodologies~Machine learning}
\ccsdesc[300]{Security and privacy~Intrusion/anomaly detection and malware mitigation}

\keywords{TRIZ, AI Agents, AI for Science, Cybersecurity}


% Mark the corresponding email visually without changing the email value.
\let\trizaOriginalNoLinkUrl\nolinkurl
\renewcommand{\nolinkurl}[1]{%
  \trizaOriginalNoLinkUrl{#1}%
  \ifstrequal{#1}{liuwy@gdut.edu.cn}{\textsuperscript{*}}{}}
\maketitle
\let\nolinkurl\trizaOriginalNoLinkUrl


\clearpage
\appendix
\counterwithin{table}{section}
\counterwithin{figure}{section}
\counterwithin{equation}{section}

\section*{Supplementary Appendix: \papertitle}
This document is the online companion to the eight-page main paper,
not an appendix appended after that paper's references. Appendices A--F
retain the public-record, mapping, contract, and result details.
Appendices G--H preserve extended method and evaluation descriptions.
Core definitions and selected results are repeated where needed for
standalone reading, not counted as new experiments or independent
replications. The motivation includes AI-assisted scientific and program
development \citep{jumper2021,romeraparedes2024,novikov2025}.

\section*{Context}


\textbf{TRIZ and computational invention.} TRIZ systematizes invention through contradictions, engineering parameters, Ideality, and reusable principles \citep{altshuller1984}; later work adapts these ideas to business and management \citep{mann2004}. Recent LLM systems automate parts of the TRIZ reasoning process or coordinate TRIZ-specialized agents \citep{jiang2025,szczepanik2025}. These systems mainly use TRIZ to structure the ideation process. \triza{} makes a narrower and testable claim: the selected principle becomes the declared type of a mutation inside an executable code-evolution loop, and that declaration is recorded before evaluation.

\textbf{LLM-guided program evolution.} FunSearch \citep{romeraparedes2024} and AlphaEvolve \citep{novikov2025} show that LLM-generated program variants can be executed, evaluated, and selected to improve program performance. OmniScientist extends automation across the research lifecycle and human--agent collaboration \citep{shao2025omniscientist}; AutoSOTA links research papers to their code repositories and couples reproduction, repair, ideation, evaluation, and supervised improvement \citep{li2026autosota}. These systems automate a broad range of research activities, whereas \triza{} focuses on how candidate transformations are specified and constrained. Rather than asking the LLM only to \emph{generate a promising program variant}, it asks the LLM to \emph{instantiate a transformation applicable to the identified contradiction}, bounded by the selected TRIZ type, frozen protocol, and evidence obligations.

\textbf{Function-conditioned selection.} FI is the negative logarithm of the fraction of configurations that achieve a specified level of function \citep{hazen2007}. Its meaning depends on the reference population and qualification threshold. Wong et al.\ discuss the roles of variation and selection in increasing FI across evolving systems \citep{wong2023}; this does not imply that each candidate is useful merely because it is rare. In our current FI calculations, the qualification criteria, reference population, sample size, and stopping rule are fixed before evaluation. The fixed reference set supports comparisons within that contract; it does not establish that rankings are insensitive to other reference sets. Confidence sequences valid across repeated interim analyses may support future experiments with data-dependent stopping \citep{howard2021}.

\section{Traceability Artifacts}
\label{app:traceability}

This appendix identifies the supplementary materials and the checks that readers
can perform. All current materials share one public entry point:
\url{https://triza.cn/triza-paper-exp-data2026/}.

\subsection{Public Materials and Reader Verification}

The collection combines the compact PowerDuck project, contradiction, strategy,
and score records with retained experiment contracts, source code, measurements,
and decision ledgers. Table~\ref{tab:registry} groups the materials by their
purpose. The PowerDuck strategy document identifies itself as a release-normalized
summary of historical records, not a file emitted under that name during the
original run.

The English--Chinese \path{README.md}, \path{RECOVERY_REPORT.md}, and
\path{label_rules.md} help readers locate measurements, distinguish comparison
references, understand how qualification and winner labels were generated, and
identify which calculations can be replayed. Browser-readable versions are
linked from the same entry page. The verification report is a guide to retained
evidence and its limits, not a substitute for raw data or independent reproduction.

The public manifest records original and published SHA-256 values separately
for path-redacted files. The checksum list verifies the current downloadable
files, including the compact PowerDuck records; the earlier compact checksum
list and the first evidence archive are retained as versioned files.
Raw datasets, embeddings, trained weights, complete chats, account configuration,
and sample-level predictions are omitted. The provided program replays retained
Campaign08 decisions without model fitting; it does not independently reproduce
every historical detector result.

\begin{table}[htbp]
\caption{Materials in the unified public supplementary collection.}
\label{tab:registry}
\centering
\begin{tabularx}{\columnwidth}{
>{\raggedright\arraybackslash}p{0.25\columnwidth}
>{\raggedright\arraybackslash}X
>{\raggedright\arraybackslash}p{0.22\columnwidth}}
\toprule
\textbf{Material} & \textbf{Content} & \textbf{Reader check} \\
\midrule
\path{project.toml}, \path{triz_context.json} & PowerDuck configuration, contradiction, recommendations, and protected conditions & Inspect the stated problem and allowed scope \\
\path{triz_strategy_plan.md}, \path{scores.jsonl} & Release-normalized transformation summary and two recorded evaluation rows & Distinguish the transformation, comparison references, and outcome \\
\path{README.md}, \path{RECOVERY_REPORT.md} & Bilingual navigation, verified findings, reproduction limits, and remaining gaps & Identify what is available and what each file supports \\
\path{label_rules.md} & Bilingual qualification, winner, policy, and reference-set definitions & Interpret controller scores and acceptance rules \\
Measurement and decision CSV files & Retained family--seed measurements, paired results, subset identities, and recounted decisions & Check reported comparisons and label conditions \\
\path{originals/}, \path{replay_campaign08.py} & Retained contracts, code, outputs, and a no-training replay entrypoint & Recompute stored Campaign08 decisions \\
\path{PUBLIC_MANIFEST.json} & File attributes and separate original/public hashes where applicable & Distinguish unchanged evidence from public derivatives \\
\path{SHA256SUMS.txt} & Digests for the current downloadable files & Check file identity and integrity \\
\bottomrule
\end{tabularx}
\end{table}

\noindent\begin{minipage}{\linewidth}
\subsection{Summary of Reported PowerDuck Results}

The following summary distinguishes the paired comparison baseline from
the incumbent immediately before the recorded transformation. Principle
labels are kept at their recorded stages rather than assigned retrospectively.


\begin{verbatim}
{
  "task": "PowerDuck GOOSE",
  "comparison": "same_protocol_paired_rerun",
  "desired_parameter": "weighted_f1",
  "undesired_parameter": "generalization_risk",
  "diagnostic_principle": "Local Quality",
  "selection_principles": ["Feedback", "Prior Action"],
  "execution_principle": "Prior Action",
  "metrics": {
    "paired_baseline_weighted_f1": 0.9694326351160863,
    "prior_incumbent_weighted_f1": 0.9960030546145862,
    "candidate_weighted_f1": 0.9981678740282784
  },
  "recorded_decision": "promoted",
  "boundary": "no_principle_level_causal_claim"
}
\end{verbatim}
\end{minipage}


\section{Full TRIZ-to-Transformation Mapping}
\label{app:fullmap}

The mapping defines transformation classes. Each concrete operator must specify input and output types, allowed arguments, prohibited changes, and required checks. Table~\ref{tab:fullmap} lists all 40 classes. The table defines the available vocabulary; the reported experiments evaluate selected instances rather than establish the effectiveness of every principle.

\begin{table}[p]
\caption{TRIZ principles and software/agent transformation classes.}
\label{tab:fullmap}
\centering
\begin{tabularx}{\linewidth}{>{\raggedright\arraybackslash}p{0.035\linewidth}>{\raggedright\arraybackslash}p{0.225\linewidth}>{\raggedright\arraybackslash}X}
\toprule
\textbf{No.} & \textbf{Principle} & \textbf{Typed transformation class} \\
\midrule
1 & Segmentation & Partition modules, data, decisions, or search regions behind a stable interface. \\
2 & Taking out & Isolate a confounder, harmful path, privileged component, or expensive stage. \\
3 & Local quality & Dispatch context-specific policy under a registered context type. \\
4 & Asymmetry & Apply different rules to unequal classes, risks, or resource directions. \\
5 & Merging & Compose compatible stages, features, models, or evaluation passes. \\
6 & Universality & Replace duplicated components with one multi-role typed service. \\
7 & Nested doll & Introduce nested scopes, hierarchical models, or encapsulated subpipelines. \\
8 & Anti-weight & Compensate load, bias, latency, or resource pressure with a countermeasure. \\
9 & Preliminary anti-action & Install a guard, rollback point, adversarial check, or preventive constraint. \\
10 & Prior action & Precompute, cache, initialize, or validate a state before the critical path. \\
11 & Beforehand cushioning & Add fallback, buffer, retry budget, or damage-limiting boundary. \\
12 & Equipotentiality & Normalize representations or equalize operating conditions. \\
13 & The other way round & Invert control, data flow, labels, optimization direction, or dependency. \\
14 & Spheroidality/\allowbreak curvature & Replace brittle discrete geometry with smooth or continuous structure. \\
15 & Dynamics & Adapt a parameter or module at runtime within frozen bounds. \\
16 & Partial/\allowbreak excessive action & Use approximation, overprovisioning, or a bounded partial result. \\
17 & Another dimension & Add a feature, hierarchy, temporal axis, view, or latent coordinate. \\
18 & Mechanical vibration & Introduce controlled probing, perturbation, or oscillatory scheduling. \\
19 & Periodic action & Re-evaluate, refresh, checkpoint, or sample on a registered schedule. \\
20 & Continuity of useful action & Stream or persist useful work while avoiding idle transitions. \\
21 & Skipping & Add an early exit, fast path, pruning rule, or bypass for safe cases. \\
22 & Blessing in disguise & Convert errors, attacks, or rejected candidates into training or test evidence. \\
23 & Feedback & Update a bounded controller from measured residuals or evaluator evidence. \\
24 & Intermediary & Insert an adapter, proxy, translator, broker, or validation layer. \\
25 & Self-service & Add self-monitoring, self-configuration, or automated maintenance. \\
26 & Copying & Use a surrogate, mock, digital twin, replay, or immutable snapshot. \\
27 & Cheap short-lived objects & Use disposable sandboxes, ephemeral workers, or low-cost probes. \\
28 & Mechanics substitution & Replace a mechanism, representation, solver, or execution substrate. \\
29 & Pneumatics/\allowbreak hydraulics & Introduce elastic resource pools, queues, or flow-control abstractions. \\
30 & Flexible shells/\allowbreak thin films & Add a lightweight interface, wrapper, boundary, or soft constraint. \\
31 & Porous materials & Use sparse structure, slots, sampling windows, or selective permeability. \\
32 & Color changes & Change encoding, annotation, visibility, or diagnostic representation. \\
33 & Homogeneity & Align schemas, units, interfaces, distributions, or component types. \\
34 & Discarding/\allowbreak recovering & Garbage-collect, reset, recycle, or restore a component after use. \\
35 & Parameter Change & Change a type, range, threshold, loss, weight, or operating point. \\
36 & Phase transitions & Switch between discrete modes or representations at a registered boundary. \\
37 & Thermal expansion & Scale capacity or spacing as load changes within resource gates. \\
38 & Strong oxidants & Increase validation strength, selection pressure, or stress testing. \\
39 & Inert atmosphere & Isolate execution in a sandbox, no-op mode, or low-interaction context. \\
40 & Composite materials & Construct an ensemble or layered architecture with typed interfaces. \\
\bottomrule
\end{tabularx}
\end{table}

\clearpage

\section{Evaluation Contracts and Experimental Details}
\label{app:contracts}

This appendix specifies the information that makes a comparison well defined: the frozen evaluation conditions, reference-population rules, qualification criteria, and statistical assumptions.

\subsection{Frozen Contract}

A valid comparison binds dataset version and license; inclusion and exclusion rules; split indices; preprocessing and feature-fit scope; candidate budget; seeds; primary and secondary metrics; useful and harmful qualification events; resource ceilings; missing-data behavior; stopping rule; and code and environment digests. The protocol fingerprint is the SHA-256 of a canonical serialization of these fields. Changing any bound field creates a new protocol rather than silently extending the old result.

At evidence level $l$, the useful and harm criteria are evaluated using
recorded predictions and execution checks. Let $g_{l,k}(s)$ be the signed
residual for gate $k$, with a positive value indicating a violation.
Promotion requires
\[
G_l(s)=\mathbf{1}\{g_{l,k}(s)\le 0\ \text{for every registered gate }k\}=1.
\]
For a minimum required improvement, the residual is the required improvement
minus the measured improvement. No FI or Ideality value can override $G_l(s)=0$.
The current FI implementations use fixed-sample calculations, not
anytime-valid confidence sequences. Repeated monitoring or data-dependent
stopping would require an appropriate sequential method
\citep{howard2021}.

The counts in Equations~\ref{eq:usefulfi} and~\ref{eq:harmfi} assign equal
weight to configurations in the fixed reference set. For multiple harm
measures, the FI V2 implementation counts a reference only if it is
at least as harmful on every included coordinate, with larger values
defined as worse. Exhaustive counts describe that finite set; confidence
claims about a larger population additionally require a stated sampling
distribution and valid sampling assumptions.

For the proposed conservative rule, a fixed-sample implementation can use
one-sided Clopper--Pearson bounds for independent Bernoulli reference
draws and a candidate fixed before those draws are inspected.
At joint coverage of at least $1-\alpha$ for one fixed candidate,
allocate $\alpha/2$ to the useful-probability upper bound and $\alpha/2$
to the harm-probability lower bound. Transform them by $-\log_2$ to obtain
$I_{U,l}^{\mathrm{LB}}$ and $I_{H,l}^{\mathrm{UB}}$.
This is a planned inference specification, not a coverage guarantee for
the reported adaptive searches or a claim that both bounds were used in
the historical campaigns. A zero harm-probability lower bound must not be
replaced by a positive pseudocount solely to obtain a finite upper bound.

\subsection{Dataset and Metric Contracts}


The formal CICIoT strict gate requires a relative macro-F1 gain of at least $3.2277\%$ and detector-FPR reduction of at least $0.005$, in addition to the integrity gates. These thresholds belong to that campaign and are not universal constants.

\subsection{Reference Population, Composition, and Comparison}
\label{app:reference-population}

The FI comparison specification requires the fields in Table~\ref{tab:reference-contract}. They describe how a reference population is to be defined and compared; they are not a new claim that complete reference manifests have been released for every historical activity. The compact PowerDuck records and the additional experiment records share one public entry point but support different checks, as described in Appendix~\ref{app:reproducibility}.

\begin{table}[htbp]
\caption{Fields required to specify an FI comparison.}
\label{tab:reference-contract}
\centering
\begin{tabularx}{\linewidth}{>{\raggedright\arraybackslash}p{0.28\linewidth}>{\raggedright\arraybackslash}X}
\toprule
\textbf{Field} & \textbf{Meaning for the comparison} \\
\midrule
Allowed configurations & The transformation grammar, parameter ranges, and resource limits that define which configurations can enter the reference space. \\
Reference identity & The included configurations and their identity or equivalence rules. The record must state whether entries are distinct configurations or repeated measurements. \\
Weights or sampling & Equal counting, a declared set of configuration weights, or a specified distribution from which reference configurations are sampled. \\
Useful qualification & The useful measurements and threshold relation used to count references that match or exceed the candidate. \\
Harm qualification & The harm measurements, their directions, and the conjunction or other comparison rule used to count references. \\
Evaluation requirements & The data, splits, preprocessing, fitting scope, seeds, measurement procedure, and evidence level shared by the candidates and references. \\
Uncertainty rules & Whether the set is enumerated or sampled, sample size and stopping rule, and the assumptions required for any probability bounds. \\
\bottomrule
\end{tabularx}
\end{table}

A repeated training seed is not automatically a new configuration. The record must identify whether seeds are measurements of one configuration or part of the configuration definition, and must state how those observations contribute to reference mass. Similarly, deduplication must follow declared identity rules rather than remove configurations selectively to change a candidate's FI.

For a single scalar useful metric $u$, with larger values indicating greater useful function, consider two candidates $a$ and $b$ evaluated against the same $\mathcal{B}_l$. If $u(a)\ge u(b)$, then
\[
\{x\in\mathcal{B}_l:u(x)\ge u(a)\}
\subseteq
\{x\in\mathcal{B}_l:u(x)\ge u(b)\}.
\]
The useful hit count for $a$ cannot exceed that for $b$. Equation~\ref{eq:usefulfi} therefore gives $I_{U,l}(a)\ge I_{U,l}(b)$ under this scalar threshold rule. Changing a common reference set can alter the values, spacing, and ties, but does not reverse this weak order. This observation applies to the point estimate with the same rule and reference measure; it does not establish the behavior of arbitrary confidence-bound estimators under adaptive search.

For multi-dimensional harm, the comparison is a partial order. Two references can have different advantages on different coordinates, so a low hit count may reflect incomparability. The useful-to-harm ratio can also change its ordering when the common reference distribution changes. Fixing the distribution establishes the conditions of the comparison, rather than invariance to all alternative distributions. Raw measurements and hard gates remain necessary to determine what the conditional FI score means for the engineering task.

\section{Extended Results}
\label{app:extended-results}

The detailed detector and controller tables are placed in Section~\ref{sec:controller-audit}, and the downstream table and figure in Section~\ref{sec:downstream}. This appendix gives the ledger scope and label definitions needed to interpret those results.

\subsection{Formal V2 Evaluation Scope}

Formal V2 covers nine held-out families and seeds 42--46, giving 45 paired
measurements under the two-million-row parent protocol. Each run allocates
1.7 million rows to fitting and 0.3 million to validation, excluding the
held-out family from both. The selected candidate is balanced-logloss RF;
RF100 baseline measurements are reused from the retained cache. Selection
used the already-observed Campaign08 nine-family qualification macro-F1,
so the 45 measurements are not 45 independent agent searches.
Section~\ref{sec:controller-audit} reports its failed acceptance thresholds.

The canonical parent-protocol hash covers the recorded configuration, not
every sample-level prediction or unrecorded software default. Recovered
checkpoints link formal V2, V17, and V18 to this parent protocol. The legacy
endpoint equals the V14 checkpoint aggregate, but its summary lacks a complete
run-level identity chain; it is not counted as an independent V14 replication.

\subsection{Campaign08 Ledger Scope}

Campaign08 contains 510 overlapping family subsets and four candidates:
\path{rf_full64_gini}, \path{rf_full64_logloss},
\path{rf_full64_balanced_logloss}, and
\path{precision_gated_rf_blend}.
Its training budget is 60,000 rows. The ledger records decision-evidence
seeds 68 and 69 and reserved qualification-label seeds 70 and 71.
The recorded \mbox{\texttt{functionalTruth}} field has 1,383 positive and 657
negative labels and is identical to the \texttt{truth} field used for the
reported scores. The separate \mbox{\texttt{promotionTruth}} field has 390
positive labels and is not the target of these metrics.
For example, balanced V2 has TP=949, FP=7, TN=650, and FN=434 against
the functional-qualification labels; specificity is
$650/(650+7)=0.98935$.
The additional evidence release includes the label generator and retained
measurement records. Let $F,R,P$ be candidate macro-F1, recall, and FPR,
averaged equally over the selected held-out-family evaluations and seeds
70/71; let $F_0,R_0,P_0$ denote the corresponding RF100 averages.
The functional-qualification label is true exactly when
\begin{equation*}
\begin{aligned}
&F\geq 0.65,\quad R\geq 0.70,\quad P\leq 0.10,\\
&(F-F_0)/F_0\geq 0.015,\quad R-R_0\geq 0,\\
&P-P_0\leq -0.005,\quad \text{and all recorded integrity gates pass}.
\end{aligned}
\end{equation*}
If $F_0=0$, the implementation sets relative gain to zero.
Integrity gates check split separation, exclusion of the held-out family
from fitting and validation, budgets, test labels, fresh-model fitting,
and the fixed prediction threshold. This label tests averaged measurements;
the two V2 policies instead require Level-1 certification on each evidence seed separately.

Within each subset, the winner is chosen by descending functional level,
macro-F1, and recall, then ascending FPR, construction steps, and candidate
ID. The winner rule does not first filter hard-gate failures and assigns
a winner even when no candidate qualifies. The separate label is
\begin{equation*}
\texttt{promotionTruth}
=\texttt{functionalTruth}\ \land\
(\texttt{candidateId}=\texttt{promotionWinner}).
\end{equation*}
Thus, 390 subsets have a positive winner label and 120 do not.
Neither label grants deployment permission.

The 510 subsets enumerate all nonempty proper subsets of the nine
held-out-family evaluations. They aggregate existing family results;
they are not 510 separately trained models that each exclude the whole
subset. The reference configurations are RF100 and the four candidates,
each with proposal weight one. Historical useful FI is the conservative
Level-1 qualification score assigned to certified candidates (zero
otherwise), rather than a separate point calculation using each candidate's
exact F1 as its threshold. Historical harmful FI uses a point estimate.
The recorded Ideality score has denominator
$\max\{0.25,I_H+\texttt{constructionSteps}/6\}$.
These implementation choices differ from the general bound-based
specification in Equation~\ref{eq:ideality}.

Original V1 requires 5\% relative evidence F1 gain and hard-gate success.
Calibrated V1 uses a 1.5085\% threshold without an additional hard-gate
test in that branch. Safety-first V2 requires new Level-1 certification,
matching evidence identities, nondecreasing positive conservative FI,
zero unresolved mass, and an Ideality comparison with the other candidates.
Balanced V2 retains certification, identity checks, positive conservative
FI, and a 3.2277\% gain threshold but removes the direct Ideality veto.
Both calibrated thresholds were developed from Campaign04, not from
Campaign08's qualification seeds. The disclosed rules and record replay
support within-corpus qualification agreement, not an external gold standard
or independent-corpus generalization.

Local timeline checks support creation of the confirmation contract and
overlay before the recorded Campaign08 launch. They are not externally
signed preregistration, and the common generator's original-time Git commit
has not been established. The public release omits private timeline excerpts.

\subsection{V14--V18 Downstream Sequence}

V17 was selected using 28 already-observed pairs and then completed the
remaining 17. It passed a separate champion rule requiring all 45 pairs,
macro-F1 above V14, a positive seed-block 95\% lower bound, and a recall
floor. Its strict balanced V2 result remained false.

V18's mechanism was tested against a V17 core evaluated on the same fitted
model. The same-fit mean deltas are macro-F1 $-0.00000248801$, MCC
$+0.00001581010$, detector FPR $+0.00003520929$, and recall
$+0.00019664105$. These differ from comparing V18 with the older cached
V17 snapshot. The retained acceptance checks reject V18; rounded snapshot
values alone are not the rejection rule. No model was retrained to prepare
the recovered evidence release.


\section{PowerDuck Transformation and Evaluation Summary}
\label{app:powerduck}

This appendix expands the recorded PowerDuck path in Table~\ref{tab:powerduck-trace}. The public bundle supports the following account:
\begin{itemize}
\item \textbf{Problem and contradiction:} improve weighted F1 while limiting
generalization risk under the frozen PowerDuck protocol.
The strategy summary records 15 classes and a ten-shot-per-attack-class
boundary; this is not a claim that the entire training set has 150 rows.
\item \textbf{Principles:} Local Quality appears in the upstream diagnostic;
the implementation record selects Feedback and Prior Action.
The executed ledger row names Prior Action. Parameter Change appears only
among the recommended principles.
\item \textbf{Transformation:} replace the scaled LDA incumbent with
\texttt{StandardScaler + LinearSVC}, using $C=1$,
balanced class weights, \texttt{dual=auto}, and random seed 42.
\item \textbf{Recorded checks:} the ledger notes the expected pipeline and
arguments, a synthetic fit/predict smoke check, seven passing unit tests,
and a clean diff check. These notes are not the underlying test logs.
\item \textbf{Outcome:} weighted F1 is $0.96943$ for the paired baseline,
$0.99600$ for the pre-promotion incumbent, and $0.99817$ for the candidate.
The candidate gains $2.87$ pp over the paired baseline and $0.22$ pp over
the incumbent. The recorded acceptance fields mark it as promoted.
\end{itemize}

The released strategy file explains how it was assembled from historical
records and lists their hashes. The additional release recovers the trace,
alignment, score ledger, baseline cache, and relevant source code.
One previously registered candidate-specification digest remains unmatched.
Raw predictions and the complete historical environment are not included. The example therefore supports inspection of the recorded
transformation and decision, not an independently reproduced result or
an isolated causal effect of a TRIZ principle.

\paragraph{Why the recorded direction is applicable.}
The released context describes Prior Action as preparing imbalance handling and preprocessing before model fitting. In the executed instance, class weighting and standardization precede fitting the linear estimator. The recorded hypothesis is that balanced LinearSVC can correct residual errors left by the scaled LDA incumbent while retaining a simple deterministic model. These statements explain the recorded choice of operation and its intended relationship to weighted-F1 and generalization risk. They do not imply that the principle uniquely specifies LinearSVC or that the classifier change alone proves an effect of TRIZ.

\paragraph{What the instance does and does not change.}
The candidate changes the permitted prediction pipeline while retaining the frozen embedding artifact and evaluation split. The recommended menu includes preprocessing and imbalance-handling options, but the executed instance is the fixed scaler--LinearSVC pipeline described above; the record does not justify adding an executed resampling step or a parameter grid to this account. The strategy document records the expected import, fit, and predict behavior, pipeline arguments, and acceptance checks. Its notes support inspection of those declared checks, not a claim that the published package contains a complete general-purpose semantic validator for all 40 principles.

\section{Reproducibility and Evidence Boundaries}
\label{app:reproducibility}

This appendix distinguishes the released files, results that can be checked against supplied records, and the additional material needed for independent reproduction.

\subsection{What the Records Support}

Appendix~\ref{app:traceability} describes the published PowerDuck files,
Appendix~\ref{app:contracts} specifies the evaluation rules, and
Appendix~\ref{app:powerduck} summarizes the PowerDuck example.
These materials support checking the stated result against the released
records but do not by themselves guarantee an exact rerun.
Reproduction also requires the data or retrieval instructions, executable
code, software environment, candidate configurations, and seeds.

The unified public collection includes the compact PowerDuck records and
the additional experiment records, which support different checks.
The latter includes the Campaign08 generator, contracts, measurements,
ledger, and a no-training replay program. All 2,040 records
and four policy confusion matrices can be checked against the retained
decision content. It also includes 45 paired measurement records for each
of formal V2, V17, and V18. These are stored detector measurements, not
sample-level predictions; their availability does not establish an
independent rerun of model fitting. The public manifest distinguishes
byte-identical evidence from path-redacted derivatives using separate
original and published hashes.

\subsection{Evidence-Specific Boundaries}

Formal CICIoT2023 V2 completed all 45 registered activities. Macro-F1
increased from $0.83256$ to $0.84839$, attack recall increased by $0.01802$,
and detector FPR decreased by $0.00168$. The candidate was not promoted
because its relative macro-F1 gain was below $3.2277\%$ and its
detector-FPR reduction was below $0.005$. The record therefore preserves
both the measured improvements and the failed promotion decision.

Campaign08 contains 2,040 nested decisions, not independent experimental
repetitions. Its F1, recall, and specificity describe agreement with
recorded functional-qualification labels, not the separate winner-selection labels.
We do not calculate significance from the decision count.
The secondary calibrated V1 policy has higher qualification F1 than balanced
V2 but lower specificity, so the audit supports a policy trade-off rather
than uniform superiority.

The V17/V18 sequence is not part of formal V2 candidate selection.
V17 passes a separate champion rule but fails strict V2 qualification.
V18 is rejected; comparison with its same-fit core decreases macro-F1
and increases FPR, while increasing MCC and recall. The older V17 cache
is retained as a historical snapshot, not substituted for this core.


\section{Extended Method and Record Specifications}
\label{app:method-details}

The method connects problem formulation to candidate generation and promotion. TRIZ guides the transformation to attempt; the LLM instantiates it; the evaluator measures benefits and harms; and the controller applies the acceptance rules. The following subsections specify this division of work, the principle-selection procedure, the role of FI, and the records needed to examine a decision.

\subsection{Framework Overview}
\label{sec:framework}

Figure~\ref{fig:architecture} shows the system as a closed invention loop. It contains four stages organized in three functional layers: (1) Problem formulation and evaluation contract (Stage 1), specifying the problem, constraints, and evaluation criteria; (2) Typed invention and evidence (Stages 2 and 3), the core of the system, selecting TRIZ principles and generating candidate solutions; (3) Closed-loop control (Stage 4), evaluating candidates and using the results to guide subsequent iterations. \textbf{Stage 1: contradiction identification.} A human-in-the-loop (HIL) brief specifies the target function and non-negotiable constraints. The contradiction extractor identifies the parameter to improve and the parameter that may worsen, together with testable criteria for the intended benefit and associated adverse effect. The metrics designer encodes these criteria in an executable evaluation protocol and freezes the protocol before candidate generation. \textbf{Stage 2: principle selection.} The Authority Provenance Knowledge Bank (AKBank) supplies source-tagged TRIZ mappings, prior outcomes, and measurement contracts. The selector retrieves the TRIZ principles recommended for the identified contradiction. \textbf{Stage 3: solution generation.} An LLM instantiates one selected principle as a code or configuration transformation. \textbf{Stage 4: evaluation and evolution.} The harness executes the candidate, after which the controller promotes it, requests revision, redirects the principle search, or stops the loop. A candidate and its audit record may be referred to an HIL operator for review and decision; any proposed high-risk action requires explicit authorization from a human operator with the required authority before execution.

\begin{figure}[t]
  \centering
  \includegraphics[width=\linewidth]{figures/triza_edit_framework.png}
  \Description{TRIZa's problem formulation, principle selection, solution generation, evaluation, and controller feedback loop.}
  \caption{System architecture of \triza{}.}
  \label{fig:architecture}
\end{figure}

These roles have separate responsibilities. AKBank supplies recommendations and their sources. The LLM produces a candidate under the selected contract. The evaluator applies the previously specified measurements, and the controller checks whether the evidence meets the promotion rules. A candidate's result cannot authorize changing its evaluation protocol after the fact.

\subsection{Theory-Guided Ideation}
\label{sec:ideation}

TRIZ adds a declared set of transformation directions to the feedback loop. A failed evaluation can motivate many edits; the selected principle identifies the kind of operation the agent will attempt before it writes the next candidate. For example, when attack recall must improve without increasing benign false positives, context-specific behavior suggests a different direction from changing one global threshold. TRIZ provides names and source mappings for such directions. These are design examples, not claims that each direction was used in every reported campaign.

We adopt TRIZ because its contradiction-based principles provide a reusable vocabulary for generating and explaining transformations \citep{altshuller1984,mann2004}. A principle does not specify a unique estimator, architecture, or parameter value. The LLM must still choose a concrete implementation, and the fixed evaluator determines whether that implementation improves the solution under the stated constraints. TRIZ and feedback therefore operate together: the former guides the proposed operation, and the latter tests it and informs the next iteration.

\paragraph{From a contradiction to an applicable transformation.}
The workflow uses the following procedure:
\begin{enumerate}
    \item \textbf{State the improvement and adverse effect.} Record the useful property to improve, the property that may worsen, and the protected data, interfaces, and evaluation settings.
    \item \textbf{Retrieve candidate principles and their sources.} Use the recorded contradiction mapping or knowledge-bank recommendation. A matrix recommendation must identify its source and matched pair; a fallback recommendation must be recorded as such.
    \item \textbf{Specify the transformation contract.} State how the proposed operation addresses the improvement--harm pair, what inputs and outputs it accepts, what may be edited, and what must remain unchanged. A recommendation is applicable only if it can be instantiated within this scope.
    \item \textbf{Declare the instance, execute it, and evaluate it.} Record the selected principle and concrete arguments before evaluating the candidate. Check the resulting edit and execution against the contract, then measure its benefits, harms, and gate outcomes under the fixed protocol.
\end{enumerate}
This procedure defines applicability and the required checks; it does not guarantee effectiveness. A transformation can be within the allowed scope and still fail to improve the measured outcomes. When the record is incomplete, the claimed coverage of the checks must be limited to the checks actually documented.

If several recommended principles are applicable, the selected principle must be recorded with a rationale connecting its proposed operation to the improvement and adverse effect. This is a documented choice within the available recommendations, not a rule claiming a unique or guaranteed selection. Historical records can be examined for that rationale; a missing rationale cannot be supplied by assigning a principle after observing the result.

\paragraph{Typed proposal record.}
\triza{} differs from untyped program evolution in how it specifies
candidate transformations: each transformation must instantiate a
selected TRIZ principle and conform to the corresponding typed contract. Given a parent artifact $s$, an identified technical contradiction $c$, a selected principle $r$, and a frozen evaluation protocol $\kappa$, we define the typed proposal record as
\begin{equation}
    \pi=(\mathrm{id}(s),c,r,\tau_r,\Delta,\kappa),
    \qquad
    s'=\tau_r(s;\Delta),
    \label{eq:proposal}
\end{equation}
where $\mathrm{id}(s)$ identifies the parent artifact, $\tau_r$ is the typed transformation operator associated with
principle $r$, and $\Delta$ contains its concrete arguments. The
instantiated transformation $\tau_r(\cdot;\Delta)$ maps the parent
artifact $s$ to the candidate artifact $s'$. The operator specifies its input and output types, what the transformation may not change, and the evidence required for evaluation. The LLM selects concrete arguments and generates the corresponding code edit, but the execution record must show that the resulting diff conforms to the declared operator. A transformation is considered TRIZ-guided only if its operator type is declared before evaluation; a favorable result alone is insufficient.

Table~\ref{tab:mapping} gives representative mappings. Each defines a class of operations rather than one parameter value or a complete solution. For example, Local Quality permits behavior specific to a declared context; it does not authorize unrelated architecture changes. The complete map in Appendix~\ref{app:fullmap} defines 40 transformation classes. The experiments concern specific instantiated transformations and do not validate every class across all tasks.

\begin{table}[t]
\caption{Representative TRIZ principles, allowed operations, and cybersecurity examples.}
    \label{tab:mapping}
    \centering
    \begin{tabularx}{\linewidth}{>{\raggedright\arraybackslash\bfseries}p{0.90in}>{\raggedright\arraybackslash}p{2.00in}>{\raggedright\arraybackslash}X}
    \toprule
    \rowcolor{trizblue!15}
    \textcolor{triznavy}{Principle} & \textcolor{triznavy}{Typed transformation} & \textcolor{triznavy}{Cybersecurity example} \\
    \midrule
    Segmentation (1) & Partition components or decisions behind a stable interface. & Per-family detectors with a shared output schema. \\
    Local Quality (3) & Dispatch context-specific behavior under a registered context. & Class-aware scaling or decision boundaries. \\
    \rowcolor{trizrow}
    Prior Action (10) & Prepare reusable state before the critical path. & Train-only feature cache bound to a split digest. \\
    Parameter Change (35) & Change a type, threshold, loss, weight, or operating point. & Balanced class weights or calibrated loss. \\
    \bottomrule
\end{tabularx}
\end{table}


If a candidate fails, the record identifies both the attempted transformation and the requirement it did not meet. This allows the next iteration to reconsider the instance or try a different principle while keeping the evaluation conditions fixed. The design provides an explicit connection between the proposed operation and the problem it is intended to address. Whether this improves search efficiency over other generation strategies requires a separate controlled comparison.

\subsection{Functional Information and FI/Ideality}
\label{sec:fi}

FI is relative to a declared reference population, not a
population-independent measure of solution quality. At evidence level $l$,
let $\mathcal{B}_l$ contain $N_l$ configurations evaluated under protocol
$\kappa$ and the same qualification rules. Write $x\succeq_U s$ if reference
$x$ matches or exceeds the useful function attained by candidate $s$.
With equal reference weights,
\begin{equation}
I_{U,l}(s)=-\log_2\!\left[
\frac{1+\sum_{x\in\mathcal{B}_l}\mathbf{1}[x\succeq_U s]}{1+N_l}
\right].
\label{eq:usefulfi}
\end{equation}
A larger value means fewer included references attain that useful level.
$N_l$ counts actual configurations, not theoretical combinations of
principles and parameters. The add-one correction keeps this point
estimate finite at zero hits; it is not a confidence guarantee.

Harm is evaluated separately. Write $x\succeq_H s$ when $x$ is at least
as harmful on every included harm measure, such as detector false-positive
rate (FPR), recall loss, leakage, or a resource violation. Then
\begin{equation}
I_{H,l}(s)=-\log_2\!\left[
\frac{1+\sum_{x\in\mathcal{B}_l}\mathbf{1}[x\succeq_H s]}{1+N_l}
\right].
\label{eq:harmfi}
\end{equation}
More harm-comparable references give a smaller $I_H$. With multiple measures,
two candidates can be incomparable; a small hit count need not mean large
absolute harm on every coordinate. Raw metrics and mandatory gates are
therefore reported alongside FI.

An evidence level states validation requirements. Certification records
the highest level actually passed:
\begin{equation}
\begin{aligned}
L_{\mathrm{cert}}(s)=\max\{\,l\mid&s\text{ passes the validation}\\
&\text{requirements and hard gates at }l\,\}.
\end{aligned}
\label{eq:certlevel}
\end{equation}
It is $-1$ if no level is certified. Evaluation at a level does not
automatically award certification at that level.

For same-level diagnosis and tie-breaking, we specify the conservative
FI-based Ideality ratio
\begin{equation}
J_l(s)=
\frac{I_{U,l}^{\mathrm{LB}}(s)}
{\max\{\epsilon,I_{H,l}^{\mathrm{UB}}(s)+\lambda_C C_{\mathrm{path}}(s)\}}.
\label{eq:ideality}
\end{equation}
Here $\epsilon>0$, $\lambda_C\geq0$, and $C_{\mathrm{path}}$ is a possible
path-cost term. If $p_U^{\mathrm{UB}}$ bounds useful-reference probability
above, $I_U^{\mathrm{LB}}=-\log_2p_U^{\mathrm{UB}}$; a harm-probability
lower bound gives $I_H^{\mathrm{UB}}=-\log_2p_H^{\mathrm{LB}}$.
A zero harm lower bound gives an infinite upper FI and zero ratio for a
finite numerator.

Equation~\ref{eq:ideality} is a method specification. Historical FI V2
uses harmful-FI point estimates; Campaign08 uses construction steps divided
by six as a cost proxy, not measured runtime or energy. Those results do
not validate the full bound-based rule. Supplementary Appendix C gives
the sampling and fixed-sample confidence assumptions.

All candidates in one comparison share the reference set, weights,
protocol, and qualification criteria. Changing these fields changes the
comparison. For a scalar useful metric, FI is monotone in the reference
hit count and is not automatically a better ranking than the metric.
With multidimensional harm and the ratio, rankings can depend on reference
composition. Mandatory integrity, safety, replication, and resource checks
cannot be compensated by a larger $J_l$. Policies may differ in the
subsequent score tests, as illustrated in Section~\ref{sec:controller-audit}.


\paragraph{Additional interpretation conditions.}
An exhaustive count describes the chosen finite reference set. A bound for
a larger population additionally needs a stated sampling distribution;
the add-one point correction does not itself establish coverage.
Changing reference membership, weights, or qualification rules defines a
new comparison, so candidates must be re-evaluated under the new contract.

With several harm measures, one reference can have lower FPR but higher
recall loss. The conjunctive harm order can therefore leave configurations
incomparable. A low hit count is not automatically large harm on every
coordinate, and equal bit units do not make two different qualification
conditions interchangeable. Raw metrics, coordinate definitions, and
mandatory gates remain part of the interpretation.

The useful-FI point estimate for one common scalar metric is monotone in
the reference hit count. Composition can alter FI values and ties; this is
not evidence of a superior ranking over the raw metric. There is no such
composition-invariance guarantee for the multidimensional harm order or
the useful-to-harm ratio. Appendix~\ref{app:contracts} specifies the
sampling and confidence assumptions for the proposed bound-based ratio.

\subsection{Traceability and Evidence State}
\label{sec:traceability}
The traceability schema specifies parent and candidate identifiers,
contradiction and principle IDs, transformation arguments, code diffs,
protocol fingerprints, metric references, FI estimates and bounds where
available, gate outcomes, and controller decisions.
Human-readable plans and execution records are stored separately.
The public record collection is available at
\url{https://triza.cn/triza-paper-exp-data2026/}; Appendix~\ref{app:traceability}
describes its materials and the checks they support. Bilingual reader guides
explain the comparison scopes, label rules, and limits of record replay.

The controller tracks a candidate together with its evaluation records. The architectural evidence scheme covers static checks, paired evaluation, registered replication, and independently locked or external evaluation. Campaign08's functionalLevels are contract-specific performance predicates (Appendix D), not a claim that all these architectural validation tiers were executed. Promotion requires the source diff, protocol, and evaluation records. A rejected candidate remains in the archive with the reason for rejection, allowing the next proposal to address the recorded failure while retaining the evaluation conditions.

The PowerDuck records identify a contradiction between weighted F1 and
generalization risk and record a train-only \texttt{StandardScaler + LinearSVC}
transformation. Local Quality appears in the upstream diagnostic;
Feedback and Prior Action appear in the implementation-selection record,
and the executed ledger row names Prior Action.
Parameter Change is among the recommended principles, not the executed
row's label. Appendix~\ref{app:powerduck} distinguishes the paired baseline
from the pre-promotion incumbent and summarizes the recorded checks.

Table~\ref{tab:powerduck-trace} makes this chain explicit. The released strategy identifies the pre-promotion incumbent as scaled linear discriminant analysis (LDA). The recorded rationale is to standardize the inputs and prepare class weighting before fitting a simple linear estimator, with the aim of correcting residual classification errors while limiting generalization risk. This explains the proposed instance of Prior Action. It does not mean that the principle uniquely determines LinearSVC, or that other estimators could not instantiate a related operation. The historical context reports no matrix hit; the recommendations and stage-specific labels are retained as recorded.

\begin{table}[t]
\caption{Recorded PowerDuck path from contradiction to executed transformation.}
\label{tab:powerduck-trace}
\centering
\begin{tabularx}{\linewidth}{>{\raggedright\arraybackslash}p{0.25\linewidth}>{\raggedright\arraybackslash}X}
\toprule
\textbf{Step} & \textbf{Content supported by the released record} \\
\midrule
Problem & Improve weighted-F1 while limiting generalization risk under the fixed split, embeddings, and evaluation procedure. \\
Recommendation & The historical context reports no matrix hit. Local Quality is recorded at diagnosis; Feedback and Prior Action at implementation selection; Prior Action in the executed ledger row. \\
Instance and rationale & Train-only standardization and class weighting precede fitting a fixed LinearSVC candidate. The recorded hypothesis is that this simple linear estimator can correct residual errors left by the scaled LDA incumbent. \\
Protected conditions & Do not change the evaluator, data split, final-test data, metric implementation, or frozen embeddings; do not expand the fixed candidate into a parameter grid after viewing results. \\
Checks and decision & The ledger notes pipeline and argument checks, a fit/predict smoke check, seven passing unit tests, and the promotion outcome. The published package does not contain the underlying test logs. \\
\bottomrule
\end{tabularx}
\end{table}

\section{Extended Evaluation Descriptions and Comparisons}
\label{app:evaluation-details}

This appendix preserves tables condensed or omitted from the eight-page
main paper. It does not add experiments. Appendix D gives the qualification,
winner, policy, and same-fit interpretation rules.

\subsection{Task Contracts and Comparison Scopes}
\label{sec:experimental-setup}

Table~\ref{tab:task-contracts} distinguishes primary metrics, protected
conditions, and acceptance checks. The sources are PowerDuck GOOSE
\citep{zemanek2022}, CICIoT2023 \citep{neto2023}, and CIC-DDoS2019
\citep{sharafaldin2019}. Macro-F1 gives labels equal weight within a run;
weighted-F1 uses test support. CICIoT family/seed averaging is a further
aggregation step, not nine-family multiclass F1. Matching dataset names
or rounded scores do not establish independent replication.

\begin{table}[htbp]
\caption{Registered task contracts.}
\label{tab:task-contracts}
\centering
\begin{tabularx}{\columnwidth}{
>{\raggedright\arraybackslash}p{0.19\columnwidth}
>{\raggedright\arraybackslash}p{0.22\columnwidth}
>{\raggedright\arraybackslash}p{0.24\columnwidth}
>{\raggedright\arraybackslash}X}
\toprule
\textbf{Task} & \textbf{Primary metric} & \textbf{Useful criterion} & \textbf{Key harm/integrity checks} \\
\midrule
PowerDuck & weighted-F1, 15 classes & paired improvement under identical split & no embedding refit; class coverage; ten shots per attack class \\
CICIoT2023 & macro-F1 across evaluated families & family/seed improvement and recall & detector FPR, attack recall, split and seed identity \\
CIC-DDoS2019 & weighted-F1 over folds & paired fold improvement & train-only feature fit, complete fold ledger \\
Campaign08 & qualification F1 and specificity & recorded functional qualification & false acceptance, nested decision dependence \\
\bottomrule
\end{tabularx}
\end{table}

\begin{table}[htbp]
\caption{Comparison scopes of the three CICIoT2023 reports.}
\label{tab:ciciot-scopes}
\centering
\begin{tabularx}{\linewidth}{>{\raggedright\arraybackslash}p{0.20\linewidth}>{\raggedright\arraybackslash}p{0.30\linewidth}>{\raggedright\arraybackslash}X}
\toprule
\textbf{Report} & \textbf{Comparison being reported} & \textbf{Interpretation and cross-report boundary} \\
\midrule
Legacy paired result & RF100 to the historical incumbent: macro-F1 $0.83256$ to $0.87490$. & The endpoint equals the retained V14 checkpoint aggregate; the legacy summary lacks a complete run-level identity chain. It is not an additional independent replication. \\
Formal V2 campaign & RF100 to a fixed balanced-logloss RF candidate: $0.83256$ to $0.84839$ over 45 family--seed pairs. & RF100 measurements are reused from the parent cache. The candidate was selected using observed Campaign08 qualification measurements and fails strict V2 acceptance. \\
Downstream sequence & V14 to V17; then a same-fit V17 core to V18, each over 45 pairs under the same parent protocol. & V17 passes a separate champion rule but fails strict V2. V18 is rejected. Historical snapshots and the same-fit mechanism comparison are distinguished. \\
\bottomrule
\end{tabularx}
\end{table}

Formal V2, V17, and V18 share the retained two-million-row parent protocol,
but use different candidates and acceptance rules. Legacy's F1 endpoint
equals the V14 aggregate without a complete original run-level identity
chain. The V14 recall snapshot cannot supply the missing legacy identity.

\subsection{Legacy Comparison References}
\label{sec:legacy}

Table~\ref{tab:powerduck-comparators} separates the three PowerDuck objects.
Promotion compares LinearSVC with the cached scaled-LDA incumbent; the
larger paired gain compares it with balanced RBF-SVC. Appendix E preserves
the recorded rationale, parameters, checks, and source-record limits.

\begin{table}[htbp]
\caption{PowerDuck objects and the comparison associated with each score.}
\label{tab:powerduck-comparators}
\centering
\begin{tabularx}{\linewidth}{>{\raggedright\arraybackslash}p{0.30\linewidth}>{\centering\arraybackslash}p{0.15\linewidth}>{\raggedright\arraybackslash}X}
\toprule
\textbf{Object} & \textbf{Weighted-F1} & \textbf{Role in interpreting the result} \\
\midrule
Balanced RBF-SVC paired baseline & $0.96943$ & Reference for the reported $+2.87$ pp gain; identified by the recovered protocol code and score cache. \\
Scaled LDA incumbent & $0.99600$ & The recorded solution immediately before the transformation; reference for the $+0.22$ pp promotion comparison. \\
StandardScaler + LinearSVC candidate & $0.99817$ & The fixed candidate. Both gains concern this score, but they use different references, listed above. \\
\bottomrule
\end{tabularx}
\end{table}

\subsection{Formal V2 and Controller Scores}
\label{sec:controller-audit}

The fixed formal candidate was chosen by the highest already-observed
Campaign08 full-family qualification macro-F1 within the FI/Ideality
evaluation workflow, not uniquely by the strict balanced gate.
The 45-pair evaluation improves F1, recall, and FPR but fails the
$3.2277\%$ relative-gain and $0.005$ FPR-reduction thresholds.
There are 43 positive pairs and positive means in all five seed blocks.

\begin{table}[htbp]
\centering
\caption{Detector results and controller agreement with recorded
functional-qualification labels. Detector deltas are relative to RF100;
controller scores are absolute values.}
\label{tab:ciciot-evidence}
\begin{tabularx}{\linewidth}{
>{\raggedright\arraybackslash}X
>{\centering\arraybackslash}p{0.16\linewidth}
>{\centering\arraybackslash}p{0.19\linewidth}
>{\centering\arraybackslash}p{0.19\linewidth}}
\toprule
\multicolumn{4}{l}{\textbf{Detector evaluation}} \\
Configuration & Macro-F1 & $\Delta$ attack recall & $\Delta$ detector FPR \\
\midrule
RF100 formal baseline & 0.83256 & -- & -- \\
Formal V2 selected & 0.84839 & $+0.01802$ & $-0.00168$ \\
\midrule
\multicolumn{4}{l}{\textbf{Controller qualification evaluation}} \\
Policy & F1 & Recall & Specificity \\
\midrule
Original V1 (fixed 5\%) & 0.508 & 0.346 & 0.963 \\
Safety-first V2 & 0.409 & 0.257 & 0.998 \\
Balanced V2 & 0.811 & 0.686 & 0.989 \\
Calibrated V1 (secondary) & 0.955 & 0.965 & 0.883 \\
\bottomrule
\end{tabularx}
\par
The higher-precision macro-F1 endpoints are $0.83255586481$ and
$0.84839357180$, giving $\Delta=0.01583770699$ (reported as $+0.01584$).
Displayed endpoints are rounded to five decimals; deltas are computed before rounding.
\end{table}

The controller scores target functional qualification, not detector
accuracy or winner truth. Four retained originals match their registered
hashes; public path-redacted copies have separate published hashes.
Replay matches all 2,040 decisions and all four policy confusion matrices,
but does not refit models or reconstruct missing predictions. The subsets
overlap, and their contract labels are program-generated within one corpus.

\subsection{Downstream Snapshots and Same-Fit Mechanism Comparison}
\label{sec:downstream}

Table~\ref{tab:downstream} and Figure~\ref{fig:downstream} retain the snapshots.
V17 was selected using 28 observed pairs before the remaining 17 were
completed; its champion acceptance and strict V2 rejection use different
rules. V18's same-fit test uses one fitted model for its core and added
mechanism. Its MCC is higher than the core but lower than the older cache;
its absolute MCC does not change with the comparison reference.
Appendix D supplies the exact deltas and rejection basis.

\begin{table}[htbp]
\caption{Historical downstream snapshots shown in Figure~\ref{fig:downstream}, rounded directly from full precision to five decimals. The older V17 cache is not the same-fit core used to test the V18 mechanism.}
\label{tab:downstream}
\centering
\begin{tabularx}{\linewidth}{>{\raggedright\arraybackslash}p{0.17\linewidth}>{\centering\arraybackslash}p{0.17\linewidth}>{\centering\arraybackslash}p{0.11\linewidth}>{\centering\arraybackslash}p{0.23\linewidth}>{\centering\arraybackslash}X}
\toprule
\textbf{Model} & \textbf{Macro-F1} & \textbf{MCC} & \textbf{Detector FPR} & \textbf{Attack recall} \\
\midrule
V1/RF100 & 0.83256 & 0.71532 & 0.05040 & 0.90528 \\
V14 & 0.87490 & 0.77467 & 0.02285 & 0.87792 \\
V17 & 0.87530 & 0.77558 & 0.02319 & 0.88095 \\
V18 & 0.87526 & 0.77553 & 0.02324 & 0.88113 \\
\bottomrule
\end{tabularx}
\end{table}

\begin{figure}[htbp]
  \centering
  \includegraphics[width=0.82\linewidth]{figures/ciciot_v14_v18_summary.pdf}
  \Description{Macro-F1, MCC, detector FPR, and attack recall for V1, V14, V17, and V18, displayed separately.}
  \caption{Downstream CICIoT2023 snapshots. Relative to V1/RF100, V17 changes macro-F1 by $+5.13\%$, MCC by $+8.42\%$, and detector FPR by $-54.0\%$, while attack recall remains $2.43$ pp below baseline. V17 passes its separate champion rule but fails strict V2 qualification. V18 is rejected. The plotted comparison with the older V17 cache is descriptive; V18's MCC is higher than that of the same-fit V17 core but lower than that of the older cached V17. This path is separate from formal V2 selection.}
  \label{fig:downstream}
\end{figure}

\subsection{Reporting and Reproduction Limits}
\label{sec:evidence-scope}
\label{sec:reproducibility-main}

These cases illustrate the workflow, not isolated component effects.
No new ablation, independent-corpus test, or reference-population
sensitivity experiment was added. Reused caches, related candidates, and
overlapping subsets are not independent confirmations. Appendices A/F
and the bilingual guides distinguish byte identity, decision replay,
and independent refitting, and identify materials still needed.

\clearpage
\bibliographystyle{ACM-Reference-Format}
\bibliography{triaz_refs}
\end{document}
\endinput
